%% ===========================================================================
\section{Values that cannot be read from the figures}
\label{app:data}
\setcounter{table}{0}
\setcounter{figure}{0}
%% ===========================================================================

Two tables are collected here because their content is not recoverable from a
figure. Table~\ref{tab:closedform} gives the closed-form coupling factors of
Sec.~\ref{sec:verification} for the four kinematics, which no figure plots;
Table~\ref{tab:ringdown} gives the first-period decrement of each of the eight
traces of Fig.~\ref{fig:ringdown}, which a reader can see but not measure to
four digits. The coupling strength is the dimensionless parameter $x$ of
\eqref{eq:kstatic}; the reported ring-down quantity is the first-period
peak decrement converted by \eqref{eq:zeta}, not a constant modal damping
ratio when $\alpha<1$.

\begin{table}[!htbp]
\centering
\small
\caption{Coupling factor from the constitutive tensors alone for the four
kinematics, with the coupling strength $e^{\mathsf T}C^{-1}e/\kappa$ they
share. The finite element rows are mode-free values of
\eqref{eq:kstatic}, plane stress, poling $x$, window off.}
\label{tab:closedform}
% table generated from the archived results
\begin{tabular}{@{}lcc@{}}
\toprule
& $x$ & $k^2$ \\
\midrule
3-D free bar & 1.3647 & 0.5771 \\
3-D uniaxial ($k_{33}^2$) & 0.5996 & 0.3748 \\
2-D plane stress, from the record & 1.3542 & \textbf{0.5752} \\
\midrule
FEM, plane stress, free, $48\times4$ & 1.3542 & \textbf{0.57523} \\
FEM, plane stress, clamped, $24\times2$ & 1.3441 & 0.57339 \\
FEM, plane stress, clamped, $48\times4$ & 1.3480 & 0.57410 \\
FEM, plane stress, clamped, $96\times4$ & 1.3498 & 0.57444 \\
\midrule
2-D plane strain, from the record & 1.0541 & 0.5132 \\
FEM, plane strain, clamped, $48 \times 4$ & 1.0471 & 0.5115 \\
\bottomrule
\end{tabular}

\end{table}

The entries in Table~\ref{tab:closedform} separate the constitutive prediction
from the effect of the clamped boundary and mesh refinement. In
Table~\ref{tab:ringdown}, the nonmonotonic dependence on $\alpha$ reflects
the change in the early-time relaxation shape at fixed $\tau$, rather than
a change in the initial voltage.

%% generated table -- do not edit (see AUTHOR_TODO.md)
\begin{table}[!htbp]
\centering
\caption{First-period equivalent damping index of the ring-down of Fig.~\ref{fig:ringdown},
converted with \eqref{eq:zeta}, for each Caputo order. The entry is a damping
ratio only at $\alpha=1$. Mode at \SI{6855.1}{kHz}, $\tau=\SI{3}{\nano\second}$
($\omega\tau=0.13$), initial voltage \SI{143.78}{\volt}, 4000 steps. The archived step-refinement comparison gives a change of
$1.7\times10^{-5}$; this does not certify all fractional orders at that tolerance.}
\label{tab:ringdown}
\begin{tabular}{@{}l r@{}}
\toprule
$\alpha$ & $\zeta_{\rm eq}$ of the first period \\
\midrule
1.0 & 0.06463 \\
0.95 & 0.07138 \\
0.9 & 0.07819 \\
0.8 & 0.09120 \\
0.7 & 0.10127 \\
0.6 & 0.10499 \\
0.5 & 0.09943 \\
0.4 & 0.08482 \\
\bottomrule
\end{tabular}
\end{table}


\FloatBarrier
